% Teil II: der durchgerechnete Weg.

\subsection*{Gegeben und gesucht}
$a,b>0$. (a) $\sqrt{ab}\le\tfrac{a+b}2$ zeigen; (b) die Fälle mit
$\sqrt{ab}=\tfrac12(a+b)$ bestimmen.

\subsection*{a) Ungleichung $\sqrt{ab}\le\frac{a+b}2$}
Für alle reellen Zahlen ist das Quadrat nichtnegativ:
\begin{equation}
  0\le\bigl(\sqrt a-\sqrt b\bigr)^2 = a-2\sqrt{ab}+b .
\end{equation}
Damit
\begin{equation}
  2\sqrt{ab}\le a+b \;\Longrightarrow\;
  \sqrt{ab}\le\frac{a+b}2 . \tag*{$\checkmark$}
\end{equation}

\subsection*{b) Gleichheitsfall}
Gleichheit in der obigen Abschätzung genau dann, wenn
$\sqrt a-\sqrt b=0$, d.h. $\sqrt a=\sqrt b$, also $a=b$.

\begin{ergebnis}
\begin{equation*}
  \sqrt{ab}\le\frac{a+b}2 \quad\text{für }a,b>0,\qquad
  \text{Gleichheit genau für } a=b .
\end{equation*}
\end{ergebnis}

\begin{kontrolle}
Gegenprobe: $a=1,b=9$: $\sqrt{9}=3$, Durchschnitt $\frac{10}2=5$, es gilt
$3\le5$ (echt); $a=b=4$: $\sqrt{16}=4$, Durchschnitt $4$, Gleichheit. Passt.
\end{kontrolle}

\begin{merksatz}
Der mittlere Wachstumsfaktor liegt nie über dem arithmetischen Mittel;
Gleichheit nur bei gleichen Faktoren -- folgt aus $(\sqrt a-\sqrt b)^2\ge0$.
\end{merksatz}
